% !TeX root = ../main.tex

\chapter{Background}\label{chapter:background}

\section{Recommendation Systems}\label{section:background:recommendation-systems}
A recommendation system is an information system that estimates which items from a set of available options are most relevant for a particular user and presents these items as suggestions~\parencite{adomaviciusTuzhilinNextGeneration}. Its purpose is to support decision making by reducing the effort required to search through many alternatives and by prioritizing those that are most likely to match the user's needs, interests, or goals.

To do so, a recommendation system uses a representation of the user and a representation of the available items. User information may include explicit preferences, prior behavior, or contextual constraints. Item information may include structured attributes, textual descriptions, or metadata. Based on these inputs, the system identifies relevant candidates, ranks them, and presents the result to the user~\parencite{luRecommenderSystems}. Recommendation is therefore not only a retrieval task, but also a relevance estimation task.

\subsection{Main Recommendation Paradigms}
Several recommendation paradigms are commonly distinguished, each relying on a different source of evidence for estimating relevance~\parencite{adomaviciusTuzhilinNextGeneration}.

\paragraph{Collaborative Filtering.}
Collaborative filtering infers preferences from interaction patterns across many users and items. It is effective when sufficient behavioral data is available, but it is less suitable when user histories are sparse or when preferences must be interpreted in a highly specific way~\parencite{suKhoshgoftaarCollaborativeFiltering}.

\paragraph{Content-Based Recommendation.}
Content-based recommendation relies on the attributes of the items themselves. It recommends items whose properties are similar to a user's known or stated preferences. Its effectiveness depends strongly on the quality and completeness of the item descriptions~\parencite{adomaviciusTuzhilinNextGeneration}.

\paragraph{Knowledge-Based Recommendation.}
Knowledge-based recommendation uses explicit constraints, domain knowledge, or decision rules to connect user requirements with item properties. This is particularly relevant in domains where decisions are complex and cannot be inferred reliably from interaction history alone~\parencite{adomaviciusTuzhilinNextGeneration}.

\paragraph{Hybrid Recommendation.}
Hybrid recommendation combines multiple paradigms in order to compensate for the limitations of any single method. In practice, such combinations are common because user needs, item representations, and available data are heterogeneous~\parencite{burkeHybridSurvey}.

\subsection{Travel Recommendation as a Special Case}
Travel recommendation is especially challenging because destination choice depends on a combination of explicit constraints and subjective preferences~\parencite{dietzDestinationCharacterization}. Users may specify budget, season, or travel duration, while also expressing qualitative expectations related to climate, activities, atmosphere, or the intended type of experience~\parencite{aleneziTravelPersonality}. These requirements are often incomplete, ambiguous, or difficult to reduce to a fixed set of attributes.

In addition, historical interaction data can be less informative in travel recommendation than in domains with denser and more repetitive feedback. Destinations are also complex items that must be described through both structured and descriptive information. As a result, travel recommendation often requires a combination of filtering, retrieval, and richer preference interpretation~\parencite{dietzDestinationCharacterization}.

\section{Large Language Models}\label{section:background:large-language-models}
Large language models are neural language processing systems trained on very large text corpora to predict likely continuations of text sequences~\parencite{zhaoSurveyLLM}. Through this objective, they learn statistical regularities of language at multiple levels, including syntax, semantic associations, discourse structure, and recurring relations between instructions and appropriate responses. In practical use, this makes them capable of generating coherent text, following prompts, and adapting their output to the immediate conversational context.

Their importance is not limited to text generation alone. Because they process input in context, large language models can derive structured interpretations from free form language. They can identify relevant entities, extract explicit constraints, relate several conditions within one utterance, resolve references to earlier statements, and reformulate vague or underspecified requests into clearer representations~\parencite{zhaoSurveyLLM}. They can also summarize longer inputs, classify intent, transform unstructured descriptions into more regular outputs, and maintain local coherence across multiple dialogue turns. These operations should not be understood as reliable human like reasoning, but they do amount to a useful form of language-based inference.

These capabilities are especially relevant in settings where users communicate through natural language rather than fixed forms or filters. In such cases, the model can help bridge the gap between the flexible way users express goals and the more explicit representations that software systems typically require. For this reason, large language models are attractive components in applications that need to interpret open-ended requests, preserve context, and produce responses that remain readable and conversational~\parencite{aliannejadiGenerativeIRInteractions}.

An important problem that occurs in large language models is hallucination, meaning that the model can produce statements that are fluent and plausible in form, but unsupported, inaccurate, or entirely fabricated with respect to the underlying facts, the provided input, or the available evidence~\parencite{zhaoSurveyLLM,gaoRagSurvey}. This problem follows from the training objective: the model is optimized to predict likely continuations of text, not to verify whether a generated statement is true in a specific application context. As a result, linguistic coherence can create an impression of reliability even when the content is wrong. Hallucinations can therefore take several forms, including fabricated facts, incorrect attributions, non existent references, or claims that go beyond what the input actually supports.

In recommendation systems, such models can in principle be used in two different ways. First, they can be used as standalone recommenders that directly transform a user request into suggested items and explanations. This offers high flexibility, but it is especially vulnerable to hallucinations, unsupported claims, and recommendations that are only loosely tied to the available data~\parencite{gaoRagSurvey}. Second, they can be used as components within a broader recommendation pipeline, where they support tasks such as request interpretation, constraint extraction, and response formulation, while retrieval and candidate selection remain tied to explicit data representations. This reduces the freedom of the model to invent unsupported content, because crucial decisions remain connected to explicit and inspectable data. More generally, this distinction points to a broader limitation of natural language output: even when text is coherent and interpretable to humans, software systems still require a computational representation of meaning that can be compared, searched, and matched systematically. One important approach to this problem is the use of embeddings.

\section{Embeddings}\label{section:background:embeddings}
Embeddings are dense numerical vectors that represent words, phrases, or longer texts in a fixed dimensional form~\parencite{mikolovDistributedRepresentations}. Instead of describing meaning through manually defined categories, an embedding model assigns each input a position in a high dimensional vector space. The basic idea is simple: texts that occur in similar contexts or express similar meanings should receive similar vector representations. As a result, semantically related items tend to lie closer to each other in vector space than unrelated ones.

The generation of embeddings depends on the model and training objective. In early distributed representation approaches, vector values are learned from co-occurrence patterns in large corpora, so that words used in similar linguistic environments obtain similar representations~\parencite{mikolovDistributedRepresentations}. More recent encoder-based methods extend this idea to whole sentences or short documents. In these models, the input text is processed by a neural encoder and mapped to a single vector that captures information about its overall meaning~\parencite{reimersSentenceBert}. Once such vectors are available, texts can be compared computationally with similarity measures such as cosine similarity, which estimates how closely two vectors point in the same direction.

Within this thesis, embeddings provide a bridge between natural language preferences and destination descriptions. A user's request may refer to ``quiet coastal towns'' while the underlying data describes ``relaxed seaside destinations'' or ``small beach cities.'' Exact keyword matching can miss such correspondences, whereas sentence level embeddings support approximate semantic retrieval over descriptive text fields~\parencite{reimersSentenceBert}. At the same time, embeddings provide only a similarity signal. They do not by themselves determine which constraints are mandatory, how conflicting preferences should be resolved, or how recommendation results should be explained. They are therefore best understood as one retrieval-oriented component within a larger system, not as a complete decision mechanism on their own. Coordinating these operations requires a higher level workflow, which motivates the agentic perspective discussed next.

\section{Agentic Artificial Intelligence}\label{section:background:agentic-ai}
A practical definition of an AI agent describes it as a system that uses a large language model to decide the control flow of an application~\parencite{chaseWhatIsAgent}. In this view, the language model is not limited to producing a final response. It can also determine which step should follow next, whether additional information is needed, whether an external tool should be invoked, or whether processing should continue iteratively. Agentic artificial intelligence therefore refers to application designs in which the language model takes an active role in coordinating behavior rather than serving only as a passive text generator.

Agentic behavior can vary in degree. At a relatively simple level, a model may act as a router that selects between predefined branches. In more advanced settings, it may repeatedly evaluate intermediate results, call tools, maintain state, revise its course of action, and continue operating until a stopping condition is reached~\parencite{chaseWhatIsAgent}. The main idea is that the model does not only generate content, it also influences how the application proceeds. This introduces a level of adaptive behavior that is difficult to achieve with fixed, purely linear prompt-response pipelines.

This is useful when a task is inherently multi step, when later actions depend on earlier results, or when the system must decide dynamically how much work is still required. Agentic designs can support decomposition of complex problems into smaller operations, selective tool use, interaction with external data sources, iterative refinement, and explicit handling of intermediate state. They are therefore attractive in settings where a system must not only answer a question, but also search, compare, verify, revise, or coordinate several operations before producing an output. At the same time, these benefits come with trade offs: agentic systems are more difficult to design, evaluate, and monitor, and additional orchestration steps create more opportunities for error propagation.

This architectural style can also be applied to recommendation scenarios. Instead of treating recommendation as a single generation step, an agentic system can separate tasks such as request interpretation, constraint extraction, retrieval, ranking, and response synthesis. This can be useful when the recommendation result is not bounded to a single query, but can instead be refined incrementally over multiple steps.

\section{Hybrid Interfaces}\label{section:background:hybrid-interfaces}
Hybrid interfaces combine conversational interaction with graphical or structured user interface elements within a single workflow~\parencite{jarvinenPromptingFuture}. Rather than forcing users to choose between chat and a conventional interface, a hybrid design treats both as complementary interaction modes. Natural language can be used to express open-ended goals, clarifications, or revisions, while structured components can externalize selections, constraints, intermediate states, and results in a more compact and inspectable form~\parencite{jannachEvaluatingCRS}.

Such an approach is relevant because conversational interaction and graphical interaction solve different problems. Natural language is effective for expressing nuanced preferences, incomplete ideas, and iterative revisions. Graphical components are better suited for comparing alternatives, exposing filters and constraints, presenting recommendation evidence, and supporting direct user control. A hybrid interface therefore aims to combine flexibility with interpretability rather than privileging one interaction style at the expense of the other~\parencite{chenPuCritiquing}.

This becomes especially useful in the context of agentic AI. If a system interprets a request, retrieves information, evaluates alternatives, and refines results over several steps, then the interaction should not be limited to a final text response alone. Users may benefit from seeing which constraints are currently active, which alternatives are under consideration, what intermediate results have been retrieved, and where they can intervene to redirect the process~\parencite{chaseWhatIsAgent}. Hybrid interfaces can support this by connecting conversational exchange with explicit visual feedback and controllable interface elements.

In this way, a hybrid interface can make an agentic system more understandable and easier to steer. It can reduce the opacity that arises when multi step processing is compressed into a single textual answer, and it can help users correct misunderstandings before they propagate through later steps. At the same time, hybrid designs introduce their own challenges: they must coordinate multiple interaction modes without overloading the user, and they must decide carefully which internal states should be exposed and which should remain hidden. When designed well, however, they provide an effective interaction layer for applications that combine conversational flexibility with structured, multi step decision support.
